
\medskip
\noindent\SGARefTarget{sga3:xi:proposition:1:9}\textbf{Proposition 1.9.} ---
\textit{Let \(f:X\to S\) be a smooth morphism of preschemes.  Then
\(\Omega^1_{X/S}\) is a locally free module of finite type on \(X\),
and its rank at a point \(x\in X\) is equal to the dimension of the
fiber \(X_s\), where \(s=f(x)\), in a neighborhood of \(x\).}

This dimension is called the \emph{relative dimension of \(X\) over
\(S\) at \(x\)}.  Notice that it is zero (when \(f\) is smooth at
\(x\)) if and only if \(f\) is étale at \(x\).  It can also be
computed directly in terms of the functor \(F\) represented by \(X\),
as follows.  Let \(\xi\) be a point of \(F\) with values in a field
\(k\), ``localized at \(x\),'' that is, a morphism
\[
  \operatorname{Spec}(k)\longrightarrow F=X
\]
whose image is \(x\).  Consider the algebra
\[
  D(k)=k[t]/(t^2)
\]
of dual numbers over \(k\), regarded as an \(S\)-prescheme, and the
map
\SGARefTarget{sga3:xi:display:1:9:dual-number-map}
\[
  F(D(k))\longrightarrow F(k)
\]
induced by the augmentation \(D(k)\to k\).  Finally, let
\(F(D(k),\xi)\) be the inverse image of \(\xi\in F(k)\) under this
map.  This set is naturally endowed with the structure of a vector
space over \(k\).  In fact, it is the vector-space dual of
\[
  \Omega^1_{X/S}\otimes_{\mathscr O_{X,x}}k,
\]
and its dimension is the relative dimension of \(X\) over \(S\) at
\(x\).

To describe the vector-space law on \(F(D(k),\xi)\) explicitly, it is
useful to introduce, more generally and as in Exposé II, for every
vector space \(V\) over \(k\), the algebra
\[
  D_k(V)=k+V,
\]
where \(V\) is an ideal of square zero, and to regard
\(F(D_k(V),\xi)\), the inverse image of \(\xi\) under
\[
  F(D_k(V))\longrightarrow F(k),
\]
as a covariant functor of \(V\) with values in
\((\mathbf{Ens})\).  It is then enough that this functor commute with
products of two factors.  This means that \(F\) carries certain
amalgamated sums of a very special kind to fiber products; compare
Exposé II.  This condition always holds when \(F\) is representable.
It follows that the sets \(F(D_k(V),\xi)\), and in particular
\(F(D(k),\xi)\), are endowed with vector-space structures over \(k\).
Thus one can define the relative dimension of \(F\) over \(X\) at
the ``point'' \(\xi\) under conditions appreciably more general than
the representability of \(F\).

In the present Exposé, the fact that certain functors to be described
explicitly are representable by preschemes smooth over \(S\) will be
used chiefly through \SGARefLink{sga3:xi:proposition:1:10}{the following result}.  It will be our technical
means of passing from constructions over the completion of the local
ring \(A\) of a point \(s\) of a noetherian scheme \(S\) to a local
ring \(A'\) over \(A\) that is étale over \(A\), and hence, in
particular, of passing from \(s\) to neighboring points.

\medskip
\noindent\SGARefTarget{sga3:xi:proposition:1:10}\textbf{Proposition 1.10.} ---
\textit{Let \(f:X\to S\) be a smooth morphism of preschemes, let
\(s\in S\), and let \(x\in X\) lie over \(s\), with \(\kappa(x)\) a
finite separable extension of \(\kappa(s)\).  Then there exists a
subscheme \(S'\subset X\), étale over \(S\), passing through
\(x\).\footnote{The French re-edition prints ``a subscheme \(S'\) of
\(S\), étale over \(S\), passing through \(x\).''  Since \(x\in X\),
and both the equivalent formulation and the construction immediately
below use \(S'\to X\) and sections of \(\mathscr O_X\), the intended
reading is \(S'\subset X\).}  Thus one can find an étale morphism
\(S'\to S\), a point \(s'\in S'\) above \(s\) whose residue-field
extension is \(\kappa(x)/\kappa(s)\), and an \(S\)-morphism
\(S'\to X\) carrying \(s'\) to \(x\).}

To construct \(S'\), simply choose sections
\(g_1,\ldots,g_n\) of \(\mathscr O_X\) on a neighborhood of \(x\)
whose images at \(x\) form a regular system of parameters for the
local ring of the fiber \(X_s\) at \(x\).  The subscheme \(S'\)
defined by the \(g_i\) is then étale over \(S\) at \(x\), and, after
shrinking \(S'\), it is étale over \(S\).

We shall use
\SGARefLink{sga3:xi:proposition:1:10}{1.10} when
\(\kappa(x)=\kappa(s)\), that is, when \(x\) is rational over
\(\kappa(s)\), or equivalently when \(x\) may be regarded as a
section of \(X_s\) over \(\operatorname{Spec}\kappa(s)\).  In that
case, \SGARefLink{sga3:xi:proposition:1:10}{1.10} has the character
of an extension theorem for sections, after an étale extension of
the base.  It takes a particularly simple form in the following
special case.

\medskip
\noindent\SGARefTarget{sga3:xi:corollary:1:11}\textbf{Corollary 1.11.} ---
\textit{Under the hypotheses of
\SGARefLink{sga3:xi:proposition:1:10}{1.10}, suppose that \(S\) is
the spectrum of a henselian local ring and that
\(\kappa(x)=\kappa(s)\).  Then there exists a section of \(X\) over
\(S\) passing through \(x\).  It is uniquely determined if \(X\) is
in fact étale over \(S\) at \(x\).}

Indeed, since \(S\) is henselian, under the hypotheses of
\SGARefLink{sga3:xi:proposition:1:10}{1.10} the prescheme \(S'\)
contains an open subscheme that is finite over \(S\) and whose fiber
at \(s\) is reduced to \(x\).  Since this open subscheme is étale
over \(S\), it is isomorphic to \(S\), which proves the assertion.
Notice that when \(S\) is the spectrum of a complete local ring,
\SGARefLink{sga3:xi:proposition:1:10}{1.10} or
\SGARefLink{sga3:xi:corollary:1:11}{1.11} is more or less equivalent
to the classical ``Hensel's lemma,'' and one sometimes refers to the
result by that name.

